\documentclass[10pt,a4paper]{amsart}
\usepackage[margin=19mm]{geometry}
\usepackage{amsmath,amssymb,mathtools}
\usepackage[scr=rsfs,cal=euler]{mathalfa}
\usepackage{xfrac}
\usepackage[unicode,hidelinks,bookmarks=false]{hyperref}
\newcommand{\ie}{i.e.\ }
\newcommand{\cf}{cf.\ }
\newcommand{\resp}{respectively\ }
\newcommand{\Subsection}{\subsection}
\providecommand{\nobreakdash}{\nobreak\hbox{-}}
\setlength{\emergencystretch}{2em}
\multlinegap0pt
\usepackage{bm}
\usepackage{bbm}
\usepackage{dsfont}
\usepackage{xspace}
\usepackage{cancel}
\usepackage{mathtools}
\usepackage{amsthm}
\makeatletter
\@ifundefined{theorem}{\newtheorem{theorem}{Theorem}}{}
\@ifundefined{lemma}{\newtheorem{lemma}[theorem]{Lemma}}{}
\makeatother
\newcommand{\FF}{\mathbf{F}}
\newcommand{\KK}{\mathbf{K}}
\newcommand{\Osh}{{\mathcal O}}                        %  Structure sheaf
\newcommand{\PP}{\mathbb{P}} % projective space
\renewcommand{\leq}{\leqslant}
\title[Effective calculation of local Weil functions via presentati]{Effective calculation of local Weil functions via presentations of Cartier divisors}
\author{Nathan Grieve}
\date{Source: arXiv:2510.18284v1 (CC BY 4.0)}
\begin{document}
\maketitle

\noindent\emph{Source and licence.} Effective calculation of local Weil functions via presentations of Cartier divisors. Version arXiv:2510.18284v1 (CC BY 4.0), distributed under the Creative Commons Attribution 4.0 International licence, as recorded in the arXiv licence metadata for this exact version and shown on the abstract page https://arxiv.org/abs/2510.18284v1. Full rights notice: licenses/arxiv-2510.18284-CC-BY-4.0.txt.\par\smallskip

\noindent\textit{Editorial scope.} Counted unit: the Theorem whose source label is \texttt{explicit:local:weil:constant} in the article (source0.tex lines 682--695, source label \texttt{explicit:local:weil:constant}), delivered together with the article's own complete original proof of it (source0.tex lines 697--830, one \texttt{proof} environment of 134 lines). No mathematics is written, repaired or re-derived: statement and proof are the article's own text line for line. Required context retained uncounted: effective:Hilb:Null:lemma:2 (lines 504--547); motivational:local:Weil:Function:Calculation (lines 675--678). Every definition, display and label of those lines is retained unchanged. Omitted from this package and not redistributed: 1--503, 548--674, 679--681, 696--696, 831--878. Nothing inside a retained excerpt is omitted or altered: every retained body is delivered exactly as the article's lines are written, so no omission receipt of any kind accompanies this package. Citations: the retained excerpts carry no citation command at all, so no bibliography entry of the article is transcribed and no reference is redistributed. Reference adaptation: none was needed and none was made; every reference inside the delivered unit resolves inside it, so no unresolved reference or citation is produced. Printed slips kept literal: no printed word, symbol, hypothesis or number of the counted statement or of its proof was altered. Layout and class: the article's own class is \texttt{amsart} with packages including fancyhdr, enumitem, amsmath, amsthm, amssymb, mathrsfs, tikz-cd, marginnote; the wrapper is the lane's pinned \texttt{amsart} editorial wrapper at 10pt on A4, which restates the theorem-like environments the retained text needs and adds the article's own macro definitions that the retained text needs (\texttt{\textbackslash FF}, \texttt{\textbackslash KK}, \texttt{\textbackslash Osh}, \texttt{\textbackslash PP}, \texttt{\textbackslash leq}), with extra wrapper packages (bm, bbm, dsfont, xspace, cancel, mathtools). The delivered numbering is the unit's own. Assets: no retained span contains a figure environment, a graphics-inclusion command, a PDF-page inclusion, a table or a TikZ picture; no image file is copied into this package, the package declares no asset dependency and no image file is redistributed with it. Licence: version arXiv:2510.18284v1 of this article is distributed under the Creative Commons Attribution 4.0 International licence, as recorded in the arXiv licence metadata for this exact version and shown on the abstract page https://arxiv.org/abs/2510.18284v1. A full rights notice is delivered at licenses/arxiv-2510.18284-CC-BY-4.0.txt. Characters: every retained excerpt is pure ASCII, so the delivered engine, XeLaTeX with its default Latin Modern text font, reports no missing character.
\begin{lemma}\label{effective:Hilb:Null:lemma:2}   Consider a standard finite affine covering $\{U_\ell\}_{\ell \in I}$, defined over a finite extension $\FF / \KK$, of an affine $\KK$-variety $U$, where 
$$
U_{\ell} = \{x \in U : h_\ell \not = 0 \} \subset U_{\FF}
$$ 
are standard affine open subsets, for regular functions 
$$h_\ell \in \Gamma(U_{\FF},\Osh_{U_{\FF}}) \text{.}$$  
Suppose that 
$$E \subset U(\overline{\KK})$$ 
is an $\FF$-bounded subset.  

Then there are $\FF$-bounded subsets 
$$E_\ell \subset U_\ell(\overline{\KK})\text{,}$$ 
which are such that 
$$E = \bigcup_{\ell \in I} E_\ell\text{.}$$  

In fact, for each $\ell \in I$ set
$$
E_\ell := \left\{x \in E : |h_\ell(x)|_v = \max_{k \in I} |h_k(x)|_v \right\} \text{.}
$$
Further,
 fix a collection of $\FF$-algebra generators, $f_1,\dots,f_N$, for $\Gamma(U_{\FF},\Osh_{U_{\FF}})$ and set 
$$f_{N+1} := 1 / h_\ell \text{.}$$  
Then
$$
\sup_{x \in E_{\ell}} \left| f_{N+1} (x) \right|_v \leq (\# I)^\delta \sup_{x \in E} \max_{k \in I} |g_k(x)|_v < \infty \text{.}
$$
and if 
$$
f \in \Gamma(U_{\ell},\Osh_{U_{\ell}})
$$
is written in the form
$$
f = q(f_1,\dots,f_{N+1}) \text{,}
$$
for 
$$
q = q(X_1,\dots,X_{N+1}) = \sum_{\mathbf{j}} b_{\mathbf{j}} X^{\mathbf{j}} \in \FF[X_1,\dots,X_{N+1}] \text{,}
$$
it follows that 
$$
\sup_{x \in E_{\ell}} | f(x)|_v \leq (\# \operatorname{supp}(q))^{\delta} |q|_v \max \left(1,\sup_{x \in E_{\ell}} \max_{j=1,\dots,N+1} |f_j(x)|_v \right)^{\deg q} \text{.}
$$
Here $|q|_v := \max_{\mathbf{j} = (j_1,\dots,j_N)} |b_{\mathbf{j}}|_v $ is the Gauss norm of $q$ with respect to $v$.
\end{lemma}

\begin{equation}\label{motivational:local:Weil:Function:Calculation}
\lambda_{\mathcal{D}}(\cdot;v) = \log |\alpha^{-1}|_v + \left( \max_{\substack{k_1 = 0,\dots,n_1 \\
\ell_2 = 0,\dots,m_2}} \min_{\substack{\ell_1 = 0,\dots,m_1 \\ k_2 = 0,\dots,n_2}} \log \left| \frac{s_{1 {k_1}} t_{2 {\ell_2}}  }{  t_{1 {\ell_1}} s_{2 {k_2}}  }(\cdot) \right|_v \right) \text{.}
\end{equation}

\begin{theorem}\label{explicit:local:weil:constant}
Let $X$ be a geometrically integral projective variety defined over $\KK$.
Let $D$ be a Cartier divisor on $X$ and defined over a finite extension $\FF / \KK$ of $\KK$.
  Fix a place $v \in M_{\KK}$ and consider two presentations
  $$
\mathcal{D}_i = (s_{iD}; L_i,\mathbf{s}_i = (s_{ii,0},\dots,s_{i,n_i}); M_i, \mathbf{t}_i = (t_{i0},\dots,t_{i,m_i})) \text{ for $i = 1,2$ }
$$
of a Cartier divisor $D$ on $X$.  Then, there exists an effectively computable constant 
$$\mathrm{O}(1) < \infty$$ 
which is such that
$$
|\lambda_{\mathcal{D}_1}(\cdot;v) - \lambda_{\mathcal{D}_2}(\cdot;v)| \leq \mathrm{O}(1) \text{.}
$$
\end{theorem}

\begin{proof}

The goal is to establish existence of an effective constant 
$$
\mathrm{O}(1) < \infty 
$$
which is such that 
$$
- \mathrm{O}(1) \leq \lambda_{\mathcal{D}_1}(\cdot;v) - \lambda_{\mathcal{D}_2}(\cdot;v) \leq \mathrm{O}(1) \text{.}
$$

Recall that
$$
\lambda_{\mathcal{D}}(\cdot;v) = \lambda_{\mathcal{D}_1 - \mathcal{D}_2}(\cdot;v) 
$$
where
$$
\mathcal{D} := \mathcal{D}_1 - \mathcal{D}_2 = (s_{1D}s_{2D}^{-1}; L_1 \otimes M_2, \mathbf{s}_1 \mathbf{t}_2; M_1 \otimes L_2, \mathbf{t}_1 \mathbf{s}_2) \text{.}
$$
Thus, in light of \eqref{motivational:local:Weil:Function:Calculation}, 
our goal amounts to establishing existence of an effective constant 
$$\mathrm{O}(1) < \infty$$ 
which is such that
\begin{equation}\label{desired:Eqn:1}
- \mathrm{O}(1) \leq \max_{\substack{k_1 = 0,\dots,n_1 \\
\ell_2 = 0,\dots,m_2}} \min_{\substack{\ell_1 = 0,\dots,m_1 \\ k_2 = 0,\dots,n_2}} \log \left| \frac{s_{1 {k_1}} t_{2 {\ell_2}}  }{  t_{1 {\ell_1}} s_{2 {k_2}}  }(\cdot) \right|_v \leq \mathrm{O}(1) \text{.}
\end{equation}

In establishing \eqref{desired:Eqn:1}, by interchanging the role of $\mathcal{D}_1$ and $\mathcal{D}_2$, it suffices to establish the right most inequality---namely to prove existence of a constant $\mathrm{O}(1) < \infty$ which is such that 
\begin{equation}\label{desired:Eqn:2}
\max_k \min_{\ell} \log \left| \frac{s_k}{t_\ell}(\cdot) \right|_v \leq \mathrm{O}(1) \text{.}
\end{equation} 
Here, in \eqref{desired:Eqn:2}, to reduce subscript notation we put 
$$s_k := s_{1 {k_1}} t_{2 {\ell_2}} \text{ for $k = 0,\dots, n_1m_2$}$$ 
and 
$$t_\ell := t_{1 {\ell_1}} s_{2 {k_2}} \text{ for $\ell = 0,\dots, m_1 n_2$.}$$

To establish the upper bound \eqref{desired:Eqn:2}, choose a closed embedding of $X$ into some projective space $\PP^N$
$$
X \hookrightarrow \PP^N_{\KK} = \PP^N_{[x_0:\dots:x_N]} \text{.}
$$
For each $i=0,\dots,N$, let 
$$U_i \subset X_{\FF}$$ 
be the affine open subset
$$
U_i := \left\{ x = [x_0:\dots : x_N] \in X_{\FF} : x_i \not = 0 \right\} \text{.}
$$
Then on the affine open subsets
$$
U_{i\ell} := \{x \in U_i : t_{\ell}(x) \not = 0 \} 
$$
the functions
$$
g_{k \ell} := s_k / t_{\ell}
$$
are regular.

On the other hand, the functions
$$
f_{ij} := \frac{x_j}{x_i} \text{,}
$$
for $j = 0,\dots,N$, and each fixed $ i = 0, \dots, N$, generate $\Gamma(U_i,\Osh_{U_i})$ as an $\FF$-algebra.

So, define the sets $E_i$, for each $i = 0,\dots,N$, by the condition that
$$
E_i := \left\{ x = [x_0:\dots:x_N] \in X(\overline{\KK}) : |x_i|_v = \max_j |x_j|_v \right\} \text{.}
$$
In particular, by construction, note that if $x \in E_i$, then
$$
\max_{j=0,\dots,N} \left| f_{ij}(x) \right|_v = 1 \text{.}
$$



It follows that each of the sets $E_i$ are $\FF$-bounded in $U_i$, for $i = 0,\dots,N$.  Thus, there exists a standard affine covering $\{U_{i \ell}\}$, of each of the sets $U_i$, and $\FF$-bounded subsets
$$
E_{i \ell} \subset U_{i\ell}
$$
which are such that 
$$
E_i = \bigcup_{i} E_{i\ell}
$$
and 
$$
\sup_{x \in E_{i \ell}} \max_k \left| g_{k \ell}(x) \right|_v < \infty \text{.}
$$


Thus, since the sets $E_{i \ell}$ cover $X(\overline{\KK})$, the existence of the desired finite and effectively computable constant $\gamma < \infty$ then follows as in Lemma \ref{effective:Hilb:Null:lemma:2}.


In more explicit terms, fix $\FF$-algebra generators
$$
f_{i0} := \frac{x_0}{x_i} , \dots, \hat{f}_{ii},\dots,f_{iN}:= \frac{x_N}{x_i}, f_{iN+1} := \frac{1}{h_{i\ell}}
$$
for each $\FF$-algebra $\Gamma(U_{i\ell},\Osh_{U_{i\ell}})$.  Here
$$
h_{i\ell} \in \Gamma(U_\ell,\Osh_{U_\ell})
$$
is a suitable regular function and the notation 
$
\hat{f}_{ii} 
$
means omit
$$
f_{ii} := \frac{x_i}{x_i} = 1
$$
in the chosen generating set.

Then, if 
$$
g_{k\ell} \in \Gamma(U_{i\ell},\Osh_{U_{i\ell}})
$$
is written in the form
$$
g_{k\ell} = q_{ik\ell}(f_{i0},\dots,\hat{f_{ii}},\dots,f_{iN},f_{iN+1}) \text{,}
$$
for
$$
q_{ik\ell} = q_{ik\ell}(X_1,\dots,X_{N+1}) \in \FF[X_1,\dots,X_{N+1}] \text{,}
$$
then 
\begin{multline*}
\sup_{x \in E_{i \ell}} \left| g_{k\ell}(x) \right|_v \leq \\ 
\left( \# \operatorname{supp}(q_{ik\ell}) \right)^{\delta} \left|q_{ik\ell} \right|_v \max\left(1,\sup_{x \in E_{\ell}} \max_{j=0,\dots, \hat{i}, \dots,N+1} |f_{ij}(x)|_v\right)^{\deg q_{ik\ell}} \text{.}
\end{multline*}

The above discussion then shows that the constant $\mathrm{O}(1)$ that arises in \eqref{desired:Eqn:2} may be taken to satisfy the inequality
\begin{multline*}
\mathrm{O}(1) \leq \\
\max_{k,\ell,i} \log^+ \left( \left( \# \operatorname{supp}(q_{ik\ell}) \right)^\delta |q_{ik\ell}|_v \max\left(1,\sup_{x \in E_{\ell}} \max_{j=0,\dots, \hat{i}, \dots,N+1} |f_{ij}(x)|_v\right)^{\deg q_{ik\ell}} \right) \text{.}
\end{multline*}
The existence and effectivity of the constant $\mathrm{O}(1)$ that arises in the conclusion of Theorem \ref{explicit:local:weil:constant} then follows too because of \eqref{desired:Eqn:1} and \eqref{motivational:local:Weil:Function:Calculation}.
\end{proof}

\end{document}
